%! Author = Nhan Huynh
%! Date = 17/01/2022

% Preamble
\documentclass{../tuda-beamer}
\usepackage{xparse}
\usepackage{natbib}

% Tikz
\usetikzlibrary{calc, shapes.multipart ,chains, arrows}
\usetikzlibrary{arrows.meta}

% Commands
\NewDocumentEnvironment{linkedlist}{o}
{

    \tikzset{
        single-linked/.style={
            rectangle split,
            rectangle split parts=2,
            draw,
            rectangle split horizontal,
            thick,
            on chain,
        },
    }

    \NewDocumentCommand{\marktail}{m}{
        \draw[shorten <= 1pt, shorten >= 1pt] (##1.north east) -- (##1.south west);
        \draw[shorten <= 1pt, shorten >= 1pt] (##1.north west) -- (##1.south east);
    }
    \NewDocumentCommand{\connect}{m m}{
        \draw[*->] let \p1=(##1.two), \p2=(##1.center) in (\x1,\y2) -- (##2);
    }
    \begin{tikzpicture}
        [
        >=stealth,
        scale=\IfNoValueTF{#1}{1}{#1},
        start chain,
        transform shape,
        ]
            }
            {
    \end{tikzpicture}
}


% Title information
\authors{Simon Hock}
\authors{Nhan Huynh}
\authors{Daniel Mangold}
\date{02. Februar 2022}

% Document
\begin{document}
    \maketitle

    \begin{frame}{Überblick}
        \tableofcontents
    \end{frame}

    \begin{frame}{ListItem [5min]}
        \begin{minipage}[t]{.475\textwidth}
            \lstinputlisting[style=Java, label={lst:listitem}]{codes/ListItem.java}
        \end{minipage}
        \hfill
        \begin{minipage}[t]{.475\textwidth}
            \lstinputlisting[style=Java, label={lst:linkedlist}]{codes/MyLinkedList.java}
        \end{minipage}

        \begin{figure}[h]
            \centering
            \begin{linkedlist}[.9]
                \node (a) {head};
                \begin{scope}[every node/.style=single-linked]
                    \node[right=of {a}] (b) {key};
                    \node[right=of {b}] (c) {key};
                    \node[right=of {c}] (d) {key};
                \end{scope}

                \node[right=of {d}] (e) {.....};

                \begin{scope}[every node/.style=single-linked]
                    \node[right=of {e}] (f) {key};
                    \node[right=of {f}] (g) {key};
                \end{scope}

                \node[thick, on chain, draw,inner sep=6pt] (h) {};
                \marktail{h}

                \draw[->] (a) -- (b);
                \draw[->] (e) -- (f);
                \connect{b}{c}
                \connect{c}{d}
                \connect{d}{e}
                \connect{f}{g}
                \connect{g}{h}
            \end{linkedlist}
            \caption{Eigene Linked List-Klasse auf Basis der Vorlesung}
            \label{fig:linkedlist-example}
        \end{figure}
    \end{frame}

    \begin{frame}[c]{Analogie - Perlenkette}
        \begin{minipage}{.475\textwidth}
            \begin{itemize}
                \item \inlinejava{MyLinkedList}: Perlenkette
                \item \inlinejava{head}: Anfang der Perlenkette
                \item \inlinejava{key}: Perle
                \item \inlinejava{next}: Verbindung zwischen Perlen
            \end{itemize}
        \end{minipage}
        \hfill
        \begin{minipage}{.475\textwidth}
            \centering
            \includegraphics[width=.8\textwidth]{graphics/pearls}
            \captionof{figure}{Perlenkette}
            \label{fig:pearls}
        \end{minipage}
    \end{frame}


    \section{Rekursion [25min]}
    \label{sec:recursion}
    \begin{frame}[c]{Rekursion}
        \begin{center}
            \enquote{Eine Rekursion ist eine Funktion, die sich selbst aufruft. Dabei wird versucht
            das Problem auf eine einfachere Variante des \textbf{gleichen} Problems rekursiv zu
            reduzieren, welches dann gelöst wird.}
        \end{center}
    \end{frame}

    \subsection{Vorgehensweise bei rekursiven Aufgaben}
    \label{subsec:recursion-approach}
    \begin{frame}[c]{Vorgehensweise bei rekursiven Aufgaben}
        \begin{enumerate}
            \item Identifiziere, wie das Problem auf eine einfachere Variante des Problems
            zerlegt werden kann und direkt gelöst werden kann

            \begin{itemize}
                \item Teillösungen der Gesamtlösung
            \end{itemize}
            \item Das kleinste Problem ist ein Rekursionsanker
            \item Überlege, wie die kleinen Probleme kombiniert werden können, um die
            ursprüngliche Problemstellung zu lösen
            \item Kombiniere die Teillösungen zu einer Gesamtlösung
        \end{enumerate}
    \end{frame}

    \begin{frame}[c]{H10 H2.2}
        \begin{itemize}
            \item Wir möchten zwei Listen miteinander verknüpfen.
            \item Die Verknüpfung ist folgendermaßen definiert:
            \begin{enumerate}
                \item Sei \(d_i\) ein Element aus der Zielliste, \(s_i\) ein Element aus der
                Quellliste.
                \item Falls das \inlinejava{BiPredicate} für \((d_i, s_i)\) \inlinejava{true}
                zurückliefert, so wird das Element vor \(d_i\) in der Zielliste eingefügt.
                \begin{enumerate}
                    \item Wiederhole 2. mit \(d_{i + 1}\) und starte bei dem zuletzt
                    betrachteten \(s_i\).
                \end{enumerate}
                \item Ansonsten wird das nächste Element von der Zielliste betrachtet \(s_{i +
                1}\) und wiederhole 2.
                \item Wir sind fertig, sobald alle Elemente der Quellliste betrachtet wurden.
                \item Falls wir keine Elemente aus der Zielliste mehr betrachten, aber noch
                Elemente aus der Quelliste betrachtet werden müssen, dann werden diese am Ende
                der Zielliste angehängt.
            \end{enumerate}
        \end{itemize}
    \end{frame}

    \begin{frame}[c]{Parameter}
        \begin{itemize}
            \item \inlinejava{MyLinkedList<U> otherList}: Quelliste
            \item \inlinejava{BiPredicate<? super T, ? super U> biPred}: Prädikat bestimmt, wo
            das Element aus der Quelliste in die Zielliste eingefügt werden soll
            \item \inlinejava{Function<? super U, ? extends T> fct}: Wandelt den Typ eines
            Elements aus der Quellliste in den Typ eines Elements der Zielliste um
            \item \inlinejava{Predicate<? super U> predU}: Prädikat bestimmt, ob ein Element aus
            der Quelliste überhaupt aufgenommen werden darf. Ansonsten wird eine
            \inlinejava{MyLinkedListException} geworfen
        \end{itemize}
        \vspace{-.5em}
        \lstinputlisting[style=Java, label={lst:h2-2-signature}]{codes/H2_2_Signature.java}
    \end{frame}

    \begin{frame}[c]{Beispiel}
        \begin{itemize}
            \item Falls die Summe der beiden Elemente im \inlinejava{biPred} ungerade ist, so
            wird das Element zur Zielliste aufgenommen
            \item Zuerst müssen wir die Position von der \(1\) bestimmen
        \end{itemize}
        \begin{figure}[h]
            \centering
            \begin{linkedlist}[.9]
                \node (a_s) at (0, 0) {head\(_{\text{QuellListe}}\)};
                \begin{scope}[every node/.style=single-linked]
                    \node[right=of {a_s}] (b_s) {1};
                    \node[right=of {b_s}] (c_s) {2};
                    \node[right=of {c_s}] (d_s) {3};
                    \node[right=of {d_s}] (e_s) {4};
                \end{scope}

                \node[thick, on chain, draw,inner sep=6pt] (f_s) {};
                \marktail{f_s}

                \draw[->] (a_s) -- (b_s);
                \connect{b_s}{c_s}
                \connect{c_s}{d_s}
                \connect{d_s}{e_s}
                \connect{e_s}{f_s}

                \node (a_d) at (0, -3) {head\(_{\text{Zielliste}}\)};
                \begin{scope}[every node/.style=single-linked]
                    \node[right=of {a_d}] (b_d) {5};
                    \node[right=of {b_d}] (c_d) {6};
                    \node[right=of {c_d}] (d_d) {7};
                \end{scope}

                \node[thick, on chain, draw,inner sep=6pt] (e_d) {};
                \marktail{e_d}

                \draw[->] (a_d) -- (b_d);
                \connect{b_d}{c_d}
                \connect{c_d}{d_d}
                \connect{d_d}{e_d}

                \node (biPred) at (0, -1) {\inlinejava{biPred(T, U)}};
                \draw[->, dashed] (biPred.east) -- (b_s);
                \draw[->, dashed] (biPred.east) -- (b_d);
                \draw[->, dashed] (biPred.east) -- (b_s);
                \draw[->, dashed] (biPred.east) -- (b_d);
                \draw[->, dashed] (biPred.east) -- (c_d);
            \end{linkedlist}
            \caption{Beispiel H2.2 - Bestimme Position von \(1\)}
            \label{fig:h10-h2.2-example}
        \end{figure}
    \end{frame}

    \begin{frame}[c]
        \begin{itemize}
            \item Füge die \(1\) hinzu.
        \end{itemize}
        \begin{figure}[h]
            \centering
            \begin{linkedlist}[.9]
                \node (a_s) at (0, 0) {head\(_{\text{QuellListe}}\)};
                \begin{scope}[every node/.style=single-linked]
                    \node[right=of {a_s}] (b_s) {1};
                    \node[right=of {b_s}] (c_s) {2};
                    \node[right=of {c_s}] (d_s) {3};
                    \node[right=of {d_s}] (e_s) {4};
                \end{scope}

                \node[thick, on chain, draw,inner sep=6pt] (f_s) {};
                \marktail{f_s}

                \draw[->] (a_s) -- (b_s);
                \connect{b_s}{c_s}
                \connect{c_s}{d_s}
                \connect{d_s}{e_s}
                \connect{e_s}{f_s}

                \node (a_d) at (0, -3) {head\(_{\text{Zielliste}}\)};
                \begin{scope}[every node/.style=single-linked]
                    \node[right=of {a_d}] (b_d) {5};
                    \node[right=of {b_d}] (c_ps) {1};
                    \node[right=of {c_ps}] (c_d) {6};
                    \node[right=of {c_d}] (d_d) {7};
                \end{scope}

                \node[thick, on chain, draw,inner sep=6pt] (e_d) {};
                \marktail{e_d}

                \draw[->] (a_d) -- (b_d);
                \connect{b_d}{c_ps}
                \connect{c_ps}{c_d}
                \connect{c_d}{d_d}
                \connect{d_d}{e_d}

                \node (biPred) at (0, -1) {\inlinejava{biPred(T, U)}};
            \end{linkedlist}
            \caption{Beispiel H2.2 - Füge \(1\) hinzu}
            \label{fig:h10-h2.2-example-1}
        \end{figure}
    \end{frame}

    \begin{frame}[c]
        \begin{itemize}
            \item Bestimme Position von \(2\).
        \end{itemize}
        \begin{figure}[h]
            \centering
            \begin{linkedlist}[.9]
                \node (a_s) at (0, 0) {head\(_{\text{QuellListe}}\)};
                \begin{scope}[every node/.style=single-linked]
                    \node[right=of {a_s}] (b_s) {1};
                    \node[right=of {b_s}] (c_s) {2};
                    \node[right=of {c_s}] (d_s) {3};
                    \node[right=of {d_s}] (e_s) {4};
                \end{scope}

                \node[thick, on chain, draw,inner sep=6pt] (f_s) {};
                \marktail{f_s}

                \draw[->] (a_s) -- (b_s);
                \connect{b_s}{c_s}
                \connect{c_s}{d_s}
                \connect{d_s}{e_s}
                \connect{e_s}{f_s}

                \node (a_d) at (0, -3) {head\(_{\text{Zielliste}}\)};
                \begin{scope}[every node/.style=single-linked]
                    \node[right=of {a_d}] (b_d) {5};
                    \node[right=of {b_d}] (c_ps) {1};
                    \node[right=of {c_ps}] (c_d) {6};
                    \node[right=of {c_d}] (d_d) {7};
                \end{scope}

                \node[thick, on chain, draw,inner sep=6pt] (e_d) {};
                \marktail{e_d}

                \draw[->] (a_d) -- (b_d);
                \connect{b_d}{c_ps}
                \connect{c_ps}{c_d}
                \connect{c_d}{d_d}
                \connect{d_d}{e_d}

                \node (biPred) at (0, -1) {\inlinejava{biPred(T, U)}};
                \draw[->, dashed] (biPred.east) -- (c_s);
                \draw[->, dashed] (biPred.east) -- (c_d);
                \draw[->, dashed] (biPred.east) -- (d_d);
            \end{linkedlist}
            \caption{Beispiel H2.2}
            \label{fig:h10-h2.2-example-2}
        \end{figure}
    \end{frame}

    \begin{frame}[c]
        \begin{itemize}
            \item Füge \(2\) hinzu.
        \end{itemize}
        \begin{figure}[h]
            \centering
            \begin{linkedlist}[.9]
                \node (a_s) at (0, 0) {head\(_{\text{QuellListe}}\)};
                \begin{scope}[every node/.style=single-linked]
                    \node[right=of {a_s}] (b_s) {1};
                    \node[right=of {b_s}] (c_s) {2};
                    \node[right=of {c_s}] (d_s) {3};
                    \node[right=of {d_s}] (e_s) {4};
                \end{scope}

                \node[thick, on chain, draw,inner sep=6pt] (f_s) {};
                \marktail{f_s}

                \draw[->] (a_s) -- (b_s);
                \connect{b_s}{c_s}
                \connect{c_s}{d_s}
                \connect{d_s}{e_s}
                \connect{e_s}{f_s}

                \node (a_d) at (0, -3) {head\(_{\text{Zielliste}}\)};
                \begin{scope}[every node/.style=single-linked]
                    \node[right=of {a_d}] (b_d) {5};
                    \node[right=of {b_d}] (c_ps) {1};
                    \node[right=of {c_ps}] (c_d) {6};
                    \node[right=of {c_d}] (d_ps) {2};
                    \node[right=of {d_ps}] (d_d) {7};
                \end{scope}

                \node[thick, on chain, draw,inner sep=6pt] (e_d) {};
                \marktail{e_d}

                \draw[->] (a_d) -- (b_d);
                \connect{b_d}{c_ps}
                \connect{c_ps}{c_d}
                \connect{c_d}{d_ps}
                \connect{d_ps}{d_d}
                \connect{d_d}{e_d}

                \node (biPred) at (0, -1) {\inlinejava{biPred(T, U)}};
            \end{linkedlist}
            \caption{Beispiel H2.2 - Füge \(2\) hinzu}
            \label{fig:h10-h2.2-example-3}
        \end{figure}
    \end{frame}

    \begin{frame}[c]
        \begin{itemize}
            \item Bestimme Position von \(3\) hinzu.
        \end{itemize}
        \begin{figure}[h]
            \centering
            \begin{linkedlist}[.9]
                \node (a_s) at (0, 0) {head\(_{\text{QuellListe}}\)};
                \begin{scope}[every node/.style=single-linked]
                    \node[right=of {a_s}] (b_s) {1};
                    \node[right=of {b_s}] (c_s) {2};
                    \node[right=of {c_s}] (d_s) {3};
                    \node[right=of {d_s}] (e_s) {4};
                \end{scope}

                \node[thick, on chain, draw,inner sep=6pt] (f_s) {};
                \marktail{f_s}

                \draw[->] (a_s) -- (b_s);
                \connect{b_s}{c_s}
                \connect{c_s}{d_s}
                \connect{d_s}{e_s}
                \connect{e_s}{f_s}

                \node (a_d) at (0, -3) {head\(_{\text{Zielliste}}\)};
                \begin{scope}[every node/.style=single-linked]
                    \node[right=of {a_d}] (b_d) {5};
                    \node[right=of {b_d}] (c_ps) {1};
                    \node[right=of {c_ps}] (c_d) {6};
                    \node[right=of {c_d}] (d_ps) {2};
                    \node[right=of {d_ps}] (d_d) {7};
                \end{scope}

                \node[thick, on chain, draw,inner sep=6pt] (e_d) {};
                \marktail{e_d}

                \draw[->] (a_d) -- (b_d);
                \connect{b_d}{c_ps}
                \connect{c_ps}{c_d}
                \connect{c_d}{d_ps}
                \connect{d_ps}{d_d}
                \connect{d_d}{e_d}

                \node (biPred) at (0, -1) {\inlinejava{biPred(T, U)}};
                \node (biPred) at (0, -1) {\inlinejava{biPred(T, U)}};
                \draw[->, dashed] (biPred.east) -- (e_d);
            \end{linkedlist}
            \caption{Beispiel H2.2 - Füge \(2\) hinzu}
            \label{fig:h10-h2.2-example-4}
        \end{figure}
    \end{frame}

    \begin{frame}[c]
        \begin{itemize}
            \item Wir sind am Ende der Zielliste \(\rightarrow\) Füge die restlichen Elemente von
            der Quellliste in die Zielliste hinzu.
            \item Wir sind nun fertig.
        \end{itemize}
        \begin{figure}[h]
            \centering
            \begin{linkedlist}[.8]
                \node (a_s) at (0, 0) {head\(_{\text{QuellListe}}\)};
                \begin{scope}[every node/.style=single-linked]
                    \node[right=of {a_s}] (b_s) {1};
                    \node[right=of {b_s}] (c_s) {2};
                    \node[right=of {c_s}] (d_s) {3};
                    \node[right=of {d_s}] (e_s) {4};
                \end{scope}

                \node[thick, on chain, draw,inner sep=6pt] (f_s) {};
                \marktail{f_s}

                \draw[->] (a_s) -- (b_s);
                \connect{b_s}{c_s}
                \connect{c_s}{d_s}
                \connect{d_s}{e_s}
                \connect{e_s}{f_s}

                \node (a_d) at (0, -3) {head\(_{\text{Zielliste}}\)};
                \begin{scope}[every node/.style=single-linked]
                    \node[right=of {a_d}] (b_d) {5};
                    \node[right=of {b_d}] (c_ps) {1};
                    \node[right=of {c_ps}] (c_d) {6};
                    \node[right=of {c_d}] (d_ps) {2};
                    \node[right=of {d_ps}] (d_d) {7};
                    \node[right=of {d_d}] (e_d) {3};
                    \node[right=of {e_d}] (f_d) {4};
                \end{scope}

                \node[thick, on chain, draw,inner sep=6pt] (g_d) {};
                \marktail{g_d}

                \draw[->] (a_d) -- (b_d);
                \connect{b_d}{c_ps}
                \connect{c_ps}{c_d}
                \connect{c_d}{d_ps}
                \connect{d_ps}{d_d}
                \connect{d_d}{e_d}
                \connect{e_d}{f_d}
                \connect{f_d}{g_d}
            \end{linkedlist}
            \caption{Beispiel H2.2 - Füge die restlichen Elemente hinzu}
            \label{fig:h10-h2.2-example-5}
        \end{figure}
    \end{frame}

    \begin{frame}{Welche Fälle müssen wir generell bei ListItem betrachten?}
        \begin{itemize}
            \item Operation wird in der Mitte ausgeführt
            \item Operation wird am Anfang ausgeführt (Sonderfall einmalig)
            \item Operation wird am Ende ausgeführt (Sonderfall einmalig)
        \end{itemize}
    \end{frame}

    \begin{frame}{Operation wird in der Mitte ausgeführt}
        \begin{enumerate}
            \item Identifiziere, wie das Problem in eine einfachere Variante des Problems
            zerlegt und direkt gelöst werden kann
        \end{enumerate}

        \begin{itemize}
            \item Wir betrachten ein Element
            \item Falls das Prädikat zutrifft, dann erstelle eine Liste und füge das Element hinzu
            \item Ansonsten mache nichts
        \end{itemize}

        \begin{figure}[h]
            \centering
            \begin{linkedlist}[.9]
                \node (biPred) at (0, 0) {\inlinejava{biPred(T, U)}};
                \node[single-linked] at (0, -2) (a) {key};

                \draw[->, dashed] (biPred.south) --
                node[left] {\inlinejava{true} oder \inlinejava{false}?} (a);

            \end{linkedlist}
            \caption{Identifiziere Problem}
            \label{fig:h10-h2.2-identify-problem}
        \end{figure}
    \end{frame}

    \begin{frame}
        \begin{enumerate}
            \setcounter{enumi}{1}
            \item Das kleinste Problem ist ein Rekursionsanker
        \end{enumerate}

        \begin{itemize}
            \item Kleinstes Problem: Quellliste ist leer \(\rightarrow\) Wir müssen nichts mehr
            machen
        \end{itemize}

        \vspace{2em}

        \begin{figure}[h]
            \centering
            \begin{linkedlist}[.9]
                \node (a_s) at (0, 0) {head\(_{\text{Quellliste}}\)};
                \node[thick, on chain, draw,inner sep=6pt, right=of {a_s}] (b_s) {};
                \marktail{b_s}

                \draw[->] (a_s) -- (b_s);

                \node (a_d) at (0, -1) {head\(_{\text{Zielliste}}\)};
                \begin{scope}[every node/.style=single-linked]
                    \node[right=of {a_d}] (b_d) {5};
                    \node[right=of {b_d}] (c_d) {6};
                    \node[right=of {c_d}] (d_d) {7};
                \end{scope}

                \node[thick, on chain, draw,inner sep=6pt] (e_d) {};
                \marktail{e_d}

                \draw[->] (a_d) -- (b_d);
                \connect{b_d}{c_d}
                \connect{c_d}{d_d}
                \connect{d_d}{e_d}
            \end{linkedlist}
            \caption{Kleinstes Problem}
            \label{fig:h10-h2.2-smallest-problem}
        \end{figure}
    \end{frame}

    \begin{frame}
        \lstinputlisting[style=Java, label={lst:recursion-anchor}]{codes/H2_2_Recursion_Anchor.java}
    \end{frame}

    \begin{frame}
        \begin{enumerate}
            \setcounter{enumi}{2}
            \item Überlege, wie die kleinen Probleme kombiniert werden können, um die
            ursprüngliche Problemstellung zu lösen
        \end{enumerate}

        \begin{figure}[h]
            \centering
            \begin{linkedlist}[.9]
                \node (a_s1) at (0, 0) {head};
                \node[single-linked, right=of {a_s1}] (b_s1) {1};
                \node[thick, on chain, draw,inner sep=6pt, right=of {b_s1}] (c_s1) {};
                \marktail{c_s1}
                \draw[->] (a_s1) -- (b_s1);
                \connect{b_s1}{c_s1}

                \node (a_s2) at (0, -2) {head};
                \node[single-linked, right=of {a_s2}] (b_s2) {6};
                \node[thick, on chain, draw,inner sep=6pt, right=of {b_s2}] (c_s2) {};
                \marktail{c_s2}
                \draw[->] (a_s2) -- (b_s2);
                \connect{b_s2}{c_s2}

                %let \p1=(##1.two), \p2=(##1.center) in (\x1,\y2) -- (##2);
                \draw[->, dashed] let  \p1=(b_s1.two), \p2=(b_s1.center) in (\x1,\y2) --
                node[right] {Kombinieren} (b_s2);
            \end{linkedlist}
            \caption{Visualisierung von Kombinieren}
            \label{fig:h10-h2.2-combine}
        \end{figure}
    \end{frame}

    \begin{frame}
        \begin{enumerate}
            \setcounter{enumi}{3}
            \item Kombiniere die Teillösungen zu einer Gesamtlösung
        \end{enumerate}

        \begin{itemize}
            \item Wir können nur Elemente von der Quellliste einfügen, sofern das Prädikat
            \inlinejava{predU} erfüllt wird
        \end{itemize}

        \lstinputlisting[style=Java, label={lst:check}]{codes/H2_2_Check.java}
    \end{frame}

    \begin{frame}
        \begin{enumerate}
            \setcounter{enumi}{3}
            \item Kombiniere die Teillösungen zu einer Gesamtlösung
        \end{enumerate}

        \begin{itemize}
            \item Am besten arbeiten wir mit dem Nachfolger, damit wir immer vorher einfügen können
            \item Einfügen erfolgt nur, wenn \inlinejava{biPred} erfüllt ist
        \end{itemize}
    \end{frame}

    \begin{frame}
        \lstinputlisting[style=Java, label={lst:insert}]{codes/H2_2_Insert.java}
    \end{frame}

    \begin{frame}
        \begin{itemize}
            \item Zum Schluss müssen wir zum nächsten Element gehen und wiederholen das ganze erneut
        \end{itemize}

        \lstinputlisting[style=Java, label={lst:next}]{codes/H2_2_Next.java}
    \end{frame}

    \begin{frame}{Was fehlt noch?}
        \begin{itemize}
            \item Spezialfälle mit erstem und letztem Element
            \item Wenn wir am Ende der Zielliste angekommen sind, dann fügen wir alle Elemente
            danach ein
            \item Der \inlinejava{else}-Teil wird nur ausgeführt, wenn \inlinejava{pDest.next ==
                null} gilt, also wenn wir am Ende der Zielliste angekommen sind
        \end{itemize}

        \lstinputlisting[style=Java, label={lst:insert-last}]{codes/H2_2_Insert_Last.java}
    \end{frame}

    \begin{frame}[allowframebreaks, c]
        \lstinputlisting[style=Java, label={lst:insert-helper}]{codes/H2_2_Helper.java}
    \end{frame}

    \begin{frame}{Spezialfall: Erstes Element}
        \begin{itemize}
            \item Momentan funktioniert die Methode nur für alle Elemente nach dem ersten Element...
            \item Wie können wir das ändern? \(\rightarrow\) Wir fügen ein
            \enquote{Dummy}-Element vor der eigentlichen Liste ein, damit das erste Element
            tatsächlich das zweite ist und somit unsere Implementierung funktioniert!
            \item Der Nachfolger vom \enquote{Dummy}-Element enthält das ggf. \enquote{neue}
            erste Element der Liste
        \end{itemize}
    \end{frame}

    \begin{frame}[allowframebreaks]
        \lstinputlisting[style=Java, label={lst:insert-main}]{codes/H2_2_Main.java}
    \end{frame}


    \section{Arbeitsphase [70min]}
    \label{sec:arbeitsphase}
    \begin{frame}[c]{Arbeitsphase}
        \begin{center}
            \textbf{\LARGE Selbstständiges Arbeiten}
        \end{center}
    \end{frame}
\end{document}
